\begin{tikzpicture}[american, transform shape]

  % -------------------------------------------------------------
  % Entradas de señal (horizontales alineadas con los terminales +)
  % -------------------------------------------------------------
  % U1 con yscale=-1 tiene U1.+ en y = 3.29 (2.8 + 0.49)
  % U2 con yscale=1 tiene U2.+ en y = -3.29 (-2.8 - 0.49)
  \node[ocirc] (in1) at (0.5, 3.29) {};
  \node[anchor=east] at (in1.west) {$e_1$};
  \node[ocirc] (in2) at (0.5, -3.29) {};
  \node[anchor=east] at (in2.west) {$e_2$};

  % -------------------------------------------------------------
  % Etapa 1: Buffers con sustitucion Z1 -> R1, Z2 -> R2
  % -------------------------------------------------------------
  \node[op amp, yscale=-1] (U1) at (3.5, 2.8) {};
  \node at (3.3, 2.8) {$U_1$};
  \draw (in1) -- (U1.+);

  \node[op amp] (U2) at (3.5, -2.8) {};
  \node at (3.3, -2.8) {$U_2$};
  \draw (in2) -- (U2.+);

  % Terminales inversores
  \draw (U1.-) -- (2.3, 2.31);
  \draw (U2.-) -- (2.3, -2.31);

  % Sustitucion Z1 = R1
  \draw (2.3, 0.9) to[american resistor, l={$R_1$}] (2.3, -0.9);
  \draw (2.3, 2.31) -- (2.3, 0.9);
  \draw (2.3, -2.31) -- (2.3, -0.9);

  % Sustitucion Z2 = R2
  \draw (2.3, 1.5) node[circ]{} -- (2.8, 1.5)
        to[american resistor, l_={$R_2$}] (4.2, 1.5)
        -- (4.7, 1.5) -- (4.7, 2.8) -- (U1.out);

  \draw (2.3, -1.5) node[circ]{} -- (2.8, -1.5)
        to[american resistor, l={$R_2$}] (4.2, -1.5)
        -- (4.7, -1.5) -- (4.7, -2.8) -- (U2.out);

  % -------------------------------------------------------------
  % Etapa 2: Sustitucion Z3 -> C3 serie R3, Z4 -> R4 paralelo C4
  % -------------------------------------------------------------
  \node[op amp] (U3) at (11.0, 0.0) {};
  \node at (10.8, 0.0) {$U_3$};

  % Rama superior: Z3(s) = C3 serie R3 hacia entrada inversora U3.-
  \draw (U1.out) -- (5.2, 2.8) node[circ]{}
        to[C, l={$C_3$}] (6.6, 2.8)
        to[american resistor, l={$R_3$}] (8.3, 2.8)
        -- (8.8, 2.8) |- (U3.-);

  % Rama inferior: Z3(s) = C3 serie R3 hacia entrada no inversora U3.+
  \draw (U2.out) -- (5.2, -2.8) node[circ]{}
        to[C, l_={$C_3$}] (6.6, -2.8)
        to[american resistor, l_={$R_3$}] (8.3, -2.8)
        -- (8.8, -2.8) |- (U3.+);

  % Realimentacion Z4(s) = R4 paralelo C4 en lazo inversor
  \draw (9.5, 0.49) node[circ]{} -- (9.5, 2.3);
  \draw (9.5, 2.3) -- (9.8, 2.3) -- (9.8, 3.0)
        to[american resistor, l={$R_4$}] (11.8, 3.0)
        -- (11.8, 2.3);
  \draw (9.8, 2.3) -- (9.8, 1.6)
        to[C, l={$C_4$}] (11.8, 1.6)
        -- (11.8, 2.3);
  \draw (11.8, 2.3) -- (12.3, 2.3) -| (U3.out);
  \draw (12.3, 0.0) node[circ]{};

  % Red a masa no inversora Z4(s) = R4 paralelo C4
  \draw (9.5, -0.49) node[circ]{} -- (9.5, -1.3) -- (10.0, -1.3);
  \draw (10.0, -1.3) to[american resistor, l_={$R_4$}] (10.0, -2.6);
  \draw (10.0, -1.3) -- (11.4, -1.3)
        to[C, l={$C_4$}] (11.4, -2.6);
  \draw (10.0, -2.6) -- (11.4, -2.6);
  \draw (10.7, -2.6) node[circ]{} -- (10.7, -2.9) node[ground]{};

  % Salida filtrada e_o3
  \draw (U3.out) -- (13.0, 0.0);
  \node[ocirc] (eo3) at (13.0, 0.0) {};
  \node[anchor=west] at (eo3.east) {$e_{o3}$};

\end{tikzpicture}
